\documentclass[12pt,a4paper]{article}
\usepackage[utf8]{inputenc}
\usepackage[T1]{fontenc}
\usepackage{mathpazo}
\usepackage[margin=2.5cm]{geometry}
\usepackage{setspace}
\doublespacing
\usepackage{parskip}

\begin{document}

\begin{center}
{\Large\bfseries The Volcanic Ritual Clock: Taphonomic Calibration of Javanese Mortuary Intervals and Their Pre-Indic Austronesian Origin}
\end{center}

\bigskip

\begin{center}
{\large Mukhlis Amien$^{1,*}$ \quad Go Frendi Gunawan$^{1}$}

\medskip

$^{1}$Universitas Bhinneka Nusantara, Malang, East Java, Indonesia\\
$^{*}$Corresponding author: amien@ubhinus.ac.id
\end{center}

\bigskip

\section*{Abstract}

The Javanese \textit{slametan} death ritual follows communal commemorations at 3, 7, 40, 100, and 1000 days after death. This numerical sequence has no source in Hindu, Buddhist, Islamic, or Christian scripture; four successive religious waves modified the verbal content while preserving the temporal structure. The intervals correspond to forensic decomposition stages in acidic volcanic soil (pH 4.5--5.5), with the terminal 1000-day mark matching bone dissolution in tropical Andosol ($p = 0.008$, permutation test). The underlying belief structure---death as gradual process, ritual efficacy paralleling bodily decay---is pan-Austronesian, shared by 30 of 137 cultures in the Pulotu database. We propose that Javanese mortuary timing was empirically calibrated to observable decomposition in volcanic landscapes, and that the same volcanism that calibrates this ritual clock also buries the archaeological evidence of the people who kept it.

\section*{Keywords}

mortuary ritual, taphonomy, slametan, Austronesian, volcanic soil, Java

\section*{Acknowledgements}

The authors thank Ahmad Suwandi and the members of the `Kruntelan' WhatsApp group, whose discussions inspired this paper. We also thank the DHARMA project (ERC) for making the Nusantaran epigraphic corpus available under CC-BY 4.0, and the Pulotu database team for open access to their cross-cultural dataset. This research was conducted with the assistance of a large language model (Claude, Anthropic) for systematic literature review, statistical computation, and corpus analysis. All hypotheses, interpretations, and scientific claims are the authors' own.

\end{document}
